Nguồn: \href{https://dantri.com.vn/giao-duc-huong-nghiep/vinh-biet-gs-phan-dinh-dieu-kinh-men-20180514225309113.htm}{Báo Dân Trí}

GS Phan Đình Diệu là nhà toán học Việt Nam đầu tiên đi du học ở Liên Xô (cũ) mà luận án TSKH được đăng thành một tuyển tập riêng của Viện Toán học mang tên V. A. Steklov của Viện Hàn lâm Khoa học Liên Xô.

GS Phan Đình Diệu, nhà toán học, khoa học máy tính của Việt Nam đã qua đời lúc 10h sáng 13/5, sau một thời gian lâm bệnh. Lễ viếng GS Phan Đình Diệu bắt đầu từ 9h30 ngày 18/5/2018, tại Nhà Tang lễ Bệnh viện TWQĐ 108, số 5 Trần Thánh Tông, Hà Nội; Lễ Truy điệu vào hồi 10h30.

Linh tính lạ kỳ: Trưa chủ Nhật ngày 13/5/2018, khi đang đọc và chỉnh sửa lại bài viết cũ của mình về GS Phan Đình Diệu, tôi nhận được tin buồn: GS Phan Đình Diệu đã từ trần vào sáng hôm đó tại Hà Nội, hưởng thọ 83 tuổi.

GS TSKH Phan Đình Diệu (12/6/1936 - 13/5/2018, quê quán Can Lộc, Hà Tĩnh) là nhà toán học, nhà khoa học máy tính của Việt Nam, người có nhiều đóng góp quan trọng xây dựng nền khoa học, giáo dục, văn hóa nước nhà .

Ông là một trong những người được ghi nhận có công đầu trong việc đào tạo và phát triển ngành tin học Việt Nam, là chuyên gia trong các lĩnh vực: Toán học kiến thiết, logic toán, lý thuyết thuật toán, ôtômat và ngôn ngữ hình thức, lý thuyết mật mã và an toàn thông tin...



GS Diệu là viện trưởng đầu tiên của Viện Khoa học Tính toán và Điều khiển (nay là Viện Công nghệ Thông tin Việt Nam), Chủ tịch sáng lập Hội Tin học Việt Nam, phó Trưởng ban thường trực Ban chỉ đạo Chương trình Quốc gia về Công nghệ Thông tin khóa I (1993-1997), ...

Trên trang cá nhân của GS Phan Dương Hiệu (con trai GS Phan Đình Diệu) ngày 27/6/2017, thông báo đã mua được trên trang sách của Hội Toán học Mỹ cuốn sách in công trình khoa học của GS Phan Đình Diệu mang tên: "Some Questions in Constructive Functional Analysis" (tạm dịch "Một số vấn đề về Giải tích hàm Kiến thiết"), được xuất bản trong "Proceedings of the Steklov Institute of Mathematics", 1974.

Khi mới ra trường làm cán bộ giảng dạy trẻ của Khoa Toán-Cơ-Tin học, trường ĐH Tổng hợp HN, tôi đã nghe tin về cuốn sách này và sau đó được nhìn thấy bản tiếng Nga trong Thư viện Khoa và Trường.

Tôi không hiểu gì mấy về nội dung cuốn sách nhưng, cũng như nhiều người khác, tôi rất trân trọng và khâm phục một thành tựu khoa học ở đỉnh cao.

Ngay lúc đó tôi cũng được biết GS Phan Đình Diệu là nhà toán học Việt Nam đầu tiên đi du học ở Liên Xô (cũ) mà luận án TSKH được đăng thành một tuyển tập riêng của Viện Toán học mang tên V. A. Steklov của Viện Hàn lâm KH Liên Xô. Tôi láng máng hiểu được rằng công trình này của GS Phan Đình Diệu là những kết quả mới mẻ, hệ thống và rất quan trọng về toán học kiến thiết.

Ngoài việc nghiên cứu và ứng dụng toán học, tin học và làm quản lý ra, GS Phan Đình Diệu còn dành nhiều thời gian cho giảng dạy, nói chuyện, viết báo về toán học, tin học, khoa học, triết học, văn hóa ...

Các hoạt động khoa học này góp phần làm cho ông nổi tiếng và được hâm mộ ở trong nước và quốc tế.

Tôi còn nhớ một bài viết rất hay của GS Phan Đình Diệu trên Tạp chí Toán học và Tuổi trẻ (số 56, tháng 10-11/1970) giới thiệu về Bài toán thứ 10 do David Hilbert đề xuất năm 1900 tại Paris.

Sau chặng đường dài 70 năm “chạy tiếp sức” của nhiều nhà toán học tài năng trên thế giới, người hoàn tất chặng đường cuối cùng năm 1970 là nhà toán học Nga 23 tuổi Iu. V. Matchiasevich (1974-), sinh viên của ĐH Tổng hợp Saint Petersburg (CHLB Nga).

Bài toán được phát biểu như sau: “Hãy chỉ ra một phương pháp mà nhờ nó, sau một số hữu hạn các phép toán có thể khẳng định rằng một phương trình Diophantine là có nghiệm nguyên hay không”.

Phương trình Diophantine là phương trình có dạng P(x, y,..., z) = 0 trong đó P (x, y,..., z) là một đa thức với hệ số nguyên của các ẩn x, y, ..., z và khi giải người ta chỉ tìm các nghiệm nguyên hoặc nghiệm nguyên không âm của nó. (Như thế phương trình Fermat chỉ là một phương trình Diophantine đặc biệt).

Tiếc thay, câu trả lời cho bài toán Hilbert số 10 lại ở dạng phủ định: Không tồn tại phương pháp chung để với mọi phương trình Diophantine cho trước, có thể khẳng định được rằng nó có nghiệm nguyên hay không.

Khi trân trọng công trình khoa học xuất sắc của GS Diệu, tôi cũng rất trân trọng tình cảm và đóng góp khoa học từ những người con của GS Phan Đình Diệu như: GS Phan Dương Hiệu, PGS Phan Thị Hà Dương, Phan Thị Quỳnh Dương.

Tôi đặc biệt quý trọng những người con, người cháu, những công dân biết trân trọng, sưu tập, gìn giữ và phát huy những di sản văn hóa, khoa học, nghệ thuật, ..., do cha ông, dân tộc mình để lại.

Đấy là những người con, người cháu, những công dân hiếu nghĩa và có văn hóa cao, biết uống nước nhớ nguồn.

Người ta thường nói: Đằng sau sự thành công của người đàn ông có bóng dáng một người phụ nữ.

Đối với GS TSKH Phan Đình Diệu, đó là phu nhân Văn Thị Xuân Hương thông minh, xinh đẹp và đảm đang, cả việc nước việc nhà.

Chị Hương là em gái của PGS Văn Như Cương, một nhà toán học, một thầy giáo toán và hiệu trưởng tài năng, độc đáo với bộ râu quý hiếm.

Hai đồng nghiệp và tôi đã dịch cuốn sách của A. D. Aczel từ tiếng Anh ra tiếng Việt "Câu chuyện hấp dẫn về Bài toán Fermat", NXBGDVN, Hà Nội-2000 (178 trang, đến nay đã được tái bản 6 lần).

GS Phan Đình Diệu là một trong những người đầu tiên tôi gửi tặng cuốn sách này. GS nói với tôi: "Cám ơn Nhung đã tặng mình cuốn sách về Định lý lớn Fermat! Khi cuốn sách vừa đến tay, mình đã đọc một mạch chừng 2-3 tiếng đồng hồ liên tục, từ trang đầu đến trang cuối. Cuốn sách hay quá”!

Vĩnh biệt Anh, GS TSKH Phan Đình Diệu kính mến!

GS.TSKH Trần Văn Nhung (14/5/2018)